% ============================================================
% CHAPTER 2: DevSecOps Fundamentals and Tooling
% (Theory chapter — adapted from Ahmed's Chapter 2 pattern,
%  focused on DevSecOps instead of DevOps)
% ============================================================

\chapter{DevSecOps Fundamentals and Tooling}

\section{Introduction}

The previous chapter presented the project context, analyzed existing infrastructure security tools, and defined the work methodology. Before proceeding to the planning and implementation phases, it is essential to establish a solid theoretical foundation.

This chapter introduces the core concepts, practices, and technologies that underpin our DevSecOps Policy-as-Code microservice. We begin with an overview of the DevOps movement, then examine how DevSecOps extends it by embedding security into every stage of the software lifecycle. We explore the key paradigms our project relies upon—Infrastructure as Code and Policy-as-Code—before presenting the specific technologies and design patterns adopted throughout development.

% ============================================================
\section{DevOps Overview}
% ============================================================

\subsection{Definition of DevOps}

DevOps is a portmanteau of ``Development'' and ``Operations.'' It designates both a cultural philosophy and a set of practices aimed at unifying software development and IT operations teams. Traditionally, these functions operated in isolation: developers wrote code and passed it to operations for deployment. This separation created bottlenecks, miscommunication, and slow release cycles~\cite{devopshandbook}.

The DevOps movement emerged to resolve these issues by encouraging collaboration, automation, and shared responsibility across the entire software delivery lifecycle—from code commit to production monitoring. At its core, DevOps is built upon several foundational principles:

\begin{itemize}
    \item \textbf{Collaboration}: Development, testing, and operations teams work as a single unit with shared objectives.
    \item \textbf{Automation}: Repetitive tasks—building, testing, deploying—are automated to reduce human error and accelerate delivery.
    \item \textbf{Continuous Feedback}: Monitoring, logging, and alerting systems provide real-time feedback that drives iterative improvement.
    \item \textbf{Shared Responsibility}: Every team member bears accountability for both the development and the operational stability of the application.
\end{itemize}

By combining people, processes, and tools, DevOps enables organizations to deliver software updates more frequently, with higher reliability, and at reduced cost~\cite{devopshandbook}.

\subsection{DevOps Lifecycle}

The DevOps lifecycle is typically represented as an infinite loop encompassing eight phases, reflecting the continuous nature of modern software delivery:

\begin{enumerate}
    \item \textbf{Plan}: Requirements gathering, goal setting, and sprint planning.
    \item \textbf{Code}: Collaborative development using version control systems such as Git.
    \item \textbf{Build}: Compilation and packaging of source code into deployable artifacts.
    \item \textbf{Test}: Automated unit, integration, and end-to-end testing to verify correctness.
    \item \textbf{Release}: Versioning and preparation of artifacts for deployment.
    \item \textbf{Deploy}: Automated delivery to staging or production environments via CI/CD pipelines.
    \item \textbf{Operate}: Running and managing the application in production.
    \item \textbf{Monitor}: Observing system health, performance metrics, and user behavior.
\end{enumerate}

% TODO: INSERT DEVOPS LIFECYCLE DIAGRAM
% Create a circular infinity-loop diagram in draw.io showing the 8 phases
% \begin{figure}[H]
%     \centering
%     \includegraphics[width=0.75\textwidth]{figures/devops-lifecycle.png}
%     \caption{DevOps lifecycle — continuous feedback loop}
%     \label{fig:devops-lifecycle}
% \end{figure}

\textit{[Figure: DevOps Lifecycle Diagram — create in draw.io]}

% ============================================================
\section{From DevOps to DevSecOps}
% ============================================================

\subsection{The Security Gap in Traditional DevOps}

While DevOps dramatically improved the speed and reliability of software delivery, it initially treated security as an afterthought. In a typical DevOps pipeline, security reviews occur at the end of the development cycle—after code has been written, built, and sometimes even deployed. This post-hoc approach creates several problems:

\begin{itemize}
    \item \textbf{Late detection}: Vulnerabilities discovered late in the pipeline are significantly more expensive to remediate. Research by the National Institute of Standards and Technology (NIST) estimates that fixing a defect in production costs 30 times more than fixing it during the design phase~\cite{nist}.
    \item \textbf{Speed versus security trade-off}: When teams deploy multiple times per day, manual security reviews become a bottleneck. Faced with pressure to deliver, teams may skip or abbreviate security checks.
    \item \textbf{Knowledge silos}: Security expertise is concentrated in specialized teams that are disconnected from daily development. Developers may lack the knowledge to write secure infrastructure code.
\end{itemize}

The growing number of high-profile breaches caused by infrastructure misconfigurations—Capital One (2019), Uber (2016), Toyota (2023)—demonstrated that speed without security is a recipe for disaster.

\subsection{DevSecOps: Security as a First-Class Citizen}

DevSecOps extends the DevOps philosophy by integrating security practices into every phase of the software development lifecycle, rather than treating it as a separate stage. The fundamental principle is \textbf{``shift-left security''}: moving security checks as early as possible—ideally to the moment code is written—so that vulnerabilities are caught before they propagate downstream.

In a DevSecOps pipeline, security is not a gate at the end; it is a continuous thread woven through every phase:

\begin{itemize}
    \item \textbf{Plan}: Threat modeling and security requirements are defined alongside functional requirements.
    \item \textbf{Code}: Static analysis tools scan code for vulnerabilities as developers write it.
    \item \textbf{Build}: Dependency scanning checks for known vulnerabilities in third-party libraries.
    \item \textbf{Test}: Security-focused tests (e.g., infrastructure policy checks) run alongside functional tests.
    \item \textbf{Deploy}: Automated policy enforcement blocks insecure deployments.
    \item \textbf{Operate}: Runtime protection monitors for anomalous behavior.
    \item \textbf{Monitor}: Continuous compliance auditing ensures ongoing adherence to security policies.
\end{itemize}

% TODO: INSERT DEVSECOPS LIFECYCLE DIAGRAM
% Create a circular diagram in draw.io showing DevSecOps lifecycle
% with security integrated at every stage (highlighted in red/orange)
% \begin{figure}[H]
%     \centering
%     \includegraphics[width=0.75\textwidth]{figures/devsecops-lifecycle.png}
%     \caption{DevSecOps lifecycle — security embedded at every stage}
%     \label{fig:devsecops-lifecycle}
% \end{figure}

\textit{[Figure: DevSecOps Lifecycle Diagram — create in draw.io]}

\subsection{The Shift-Left Security Concept}

The term ``shift-left'' refers to moving an activity earlier (leftward) in the development timeline. Applied to security, it means detecting and fixing vulnerabilities during development rather than after deployment.

Figure~\ref{fig:shift-left-concept} illustrates the shift-left concept: traditional security testing occurs at the right end of the pipeline (post-deployment), while shift-left security moves checks to the coding and CI/CD stages. The earlier a vulnerability is detected, the cheaper and faster it is to fix.

% TODO: INSERT SHIFT-LEFT DIAGRAM
% Create a horizontal timeline diagram in draw.io:
% Left side: Code → Build → Test → Deploy → Production
% Show "Traditional Security" arrow pointing to the right (late)
% Show "Shift-Left Security" arrow pointing to the left (early)
% Show cost multiplier: 1x (code) → 6x (test) → 15x (deploy) → 30x (production)
% \begin{figure}[H]
%     \centering
%     \includegraphics[width=0.85\textwidth]{figures/shift-left.png}
%     \caption{Shift-left security — detecting vulnerabilities earlier reduces remediation cost}
%     \label{fig:shift-left-concept}
% \end{figure}

\textit{[Figure: Shift-Left Security Diagram — create in draw.io]}

\noindent Our DevSecOps Policy-as-Code microservice operates precisely at the shift-left position: it scans infrastructure code \textit{before} deployment, directly within the CI/CD pipeline. When a developer pushes insecure Terraform, Dockerfile, or Kubernetes code to Git, the service analyzes it within 46 milliseconds and blocks the deployment if critical violations are found.

% ============================================================
\section{Key DevSecOps Concepts}
% ============================================================

\subsection{Continuous Integration and Continuous Deployment (CI/CD)}

CI/CD is the backbone of modern DevOps and DevSecOps pipelines. \textbf{Continuous Integration (CI)} is the practice of automatically building and testing code every time a developer commits changes to a shared repository. \textbf{Continuous Deployment (CD)} extends this by automatically deploying validated code to production environments.

In the context of our project, CI/CD serves as the \textit{delivery mechanism} for security enforcement. When a developer pushes infrastructure code to Gitea (our local Git server), Jenkins detects the change and triggers a pipeline that includes a mandatory security scan via our policy service. The pipeline proceeds to the build and deploy stages only if the scan returns an ALLOW verdict. This tight integration ensures that no insecure infrastructure code reaches production.

\begin{table}[H]
\centering
\caption{Comparison of CI/CD tools}
\label{tab:cicd-comparison}
\begin{tabular}{|p{2.5cm}|p{3cm}|p{3cm}|p{3cm}|}
\hline
\textbf{Criteria} & \textbf{Jenkins} & \textbf{GitLab CI} & \textbf{GitHub Actions} \\
\hline
License & Open Source & Free/Premium & Free/Premium \\
\hline
Self-hosted & \cellcolor{green!15}\textbf{Yes} & Yes & No (cloud) \\
\hline
Pipeline as Code & Jenkinsfile & .gitlab-ci.yml & .github/workflows \\
\hline
Plugin ecosystem & \cellcolor{green!15}\textbf{1,800+ plugins} & Built-in & Marketplace \\
\hline
Offline capable & \cellcolor{green!15}\textbf{Yes} & Partial & No \\
\hline
\end{tabular}
\end{table}

\noindent\textbf{Our choice: Jenkins.} We selected Jenkins because our project demonstrates an entirely offline CI/CD workflow. Jenkins runs locally on \texttt{localhost:8080}, connects to a local Gitea server, and invokes our policy service—all without internet access. This self-contained setup proves that DevSecOps can be adopted even in air-gapped or resource-constrained environments.

\subsection{Containerization}

Containerization is the process of packaging an application together with its dependencies, libraries, and configuration files into a single, lightweight, portable unit called a \textit{container}. Unlike virtual machines, containers share the host operating system's kernel, making them significantly lighter and faster to start.

In our project, Docker serves two purposes:
\begin{enumerate}
    \item \textbf{Deployment target}: Our policy service is packaged as a Docker image, enabling one-command deployment via \texttt{docker-compose up}.
    \item \textbf{Scan target}: Our service scans Dockerfiles for security best practices (e.g., avoiding the \texttt{latest} tag, not running as root, including health checks).
\end{enumerate}

\subsection{Infrastructure as Code (IaC)}

Infrastructure as Code (IaC) is the practice of defining and managing computing infrastructure—servers, networks, storage, security groups—through machine-readable configuration files rather than manual processes. IaC enables teams to version-control their infrastructure, reproduce environments reliably, and apply the same software engineering practices (code review, testing, CI/CD) to infrastructure management~\cite{terraform}.

The three major IaC formats relevant to our project are:

\begin{table}[H]
\centering
\caption{IaC formats supported by our service}
\label{tab:iac-formats}
\begin{tabular}{|p{2.5cm}|p{4cm}|p{6cm}|}
\hline
\textbf{Format} & \textbf{Tool / Language} & \textbf{Purpose} \\
\hline
Terraform HCL & HashiCorp Terraform~\cite{terraform} & Provisioning cloud resources (AWS S3, EC2, RDS, IAM, VPC, etc.) \\
\hline
Dockerfile & Docker~\cite{docker} & Defining container images (base image, dependencies, runtime config) \\
\hline
Kubernetes YAML & Kubernetes~\cite{kubernetes} & Orchestrating containerized workloads (Pods, Deployments, Services) \\
\hline
\end{tabular}
\end{table}

\noindent While IaC brings enormous benefits, it also introduces new security risks. A misconfigured Terraform resource can expose an entire database to the internet. A Dockerfile running as root grants attackers full container access. A Kubernetes Pod without security contexts can escalate privileges across the cluster. These are precisely the risks that our policy service is designed to detect and prevent.

\subsection{Policy-as-Code}

Policy-as-Code is the practice of expressing organizational security, compliance, and operational policies as executable code rather than human-readable documents. Instead of a written guideline stating ``S3 buckets must not be publicly accessible,'' a Policy-as-Code implementation encodes this rule as a programmatic check that can be automatically enforced at machine speed~\cite{opa}.

This paradigm offers critical advantages:

\begin{itemize}
    \item \textbf{Automation}: Policies are evaluated automatically, eliminating the inconsistency and latency of manual reviews.
    \item \textbf{Consistency}: Every scan applies the same rules identically, regardless of the reviewer or time of day.
    \item \textbf{Scalability}: A single policy service can evaluate thousands of scans per day without human intervention.
    \item \textbf{Auditability}: Policy definitions are version-controlled, providing a complete audit trail of what rules were enforced and when.
    \item \textbf{Education}: When combined with detailed violation messages (as in our service), Policy-as-Code transforms enforcement into a learning opportunity for developers.
\end{itemize}

\noindent Our microservice implements the Policy-as-Code paradigm through 100 Python-based security policies, each implemented as a class with a \texttt{check()} method. This approach provides the flexibility of a general-purpose programming language (Python) with the structure and discipline of a formal policy framework.

% ============================================================
\section{Design Patterns Adopted}
% ============================================================

The architecture of our service relies on several well-established software design patterns. Each pattern was chosen to address a specific architectural challenge.

\subsection{Layered Architecture Pattern}

The Layered Architecture pattern organizes the system into horizontal layers, where each layer has a specific responsibility and communicates only with its adjacent layers. This separation of concerns ensures that changes in one layer (e.g., adding a new parser) do not affect other layers (e.g., the policy engine).

Our service implements four layers:
\begin{enumerate}
    \item \textbf{API Layer} (FastAPI) — Receives HTTP requests, validates input, returns JSON responses.
    \item \textbf{Parser Layer} — Converts raw code strings into structured data.
    \item \textbf{Normalizer Layer} — Transforms heterogeneous parsed data into a unified Resource model.
    \item \textbf{Policy Engine Layer} — Evaluates normalized resources against 100 security policies and produces ALLOW/BLOCK decisions.
\end{enumerate}

\subsection{Strategy Pattern}

The Strategy pattern defines a family of algorithms, encapsulates each one, and makes them interchangeable. In our service, the three parsers (TerraformParser, DockerfileParser, KubernetesParser) are strategies that implement the same interface. The \texttt{get\_parser(code\_type)} factory function selects the appropriate parser at runtime based on the input format.

This pattern enables \textbf{extensibility}: adding support for a new IaC format (e.g., Pulumi, AWS CDK) requires only implementing a new parser class and registering it in the factory—no existing code needs to change.

\subsection{Template Method Pattern}

The Template Method pattern defines the skeleton of an algorithm in a base class, deferring some steps to subclasses. In our service, the \texttt{BasePolicy} abstract class defines the template:

\begin{lstlisting}[language=Python, caption={BasePolicy template method}, label={lst:basepolicy}]
class BasePolicy:
    rule_id: str
    severity: str  # CRITICAL, HIGH, MEDIUM, LOW
    description: str

    def check(self, resource: Resource) -> List[Violation]:
        """Subclasses implement this method."""
        raise NotImplementedError
\end{lstlisting}

Each of the 100 security policies inherits from \texttt{BasePolicy} and implements its own \texttt{check()} method. This pattern enabled us to scale from 12 initial policies to 100 without modifying the engine—each new policy is a self-contained class.

\subsection{Factory Pattern}

The Factory pattern centralizes object creation logic. In our service, the \texttt{get\_parser()} function acts as a factory, returning the appropriate parser instance based on the \texttt{code\_type} parameter:

\begin{lstlisting}[language=Python, caption={Parser factory function}, label={lst:factory}]
def get_parser(code_type: str):
    parsers = {
        "terraform": TerraformParser(),
        "dockerfile": DockerfileParser(),
        "yaml": KubernetesParser(),
    }
    return parsers.get(code_type)
\end{lstlisting}

This pattern decouples the API layer from the specific parser implementations, allowing new formats to be added without modifying the routing logic.

% ============================================================
\section{Selected Technologies Overview}
% ============================================================

This section presents the key technologies used in our project, explaining why each was selected and how it contributes to the overall system.

\subsection{Python 3.10}

Python was selected as the primary development language for several reasons: its readability promotes maintainable code (critical for a project with 100 policy classes), its rich ecosystem includes specialized libraries for parsing HCL and YAML formats, and its support for abstract classes and dataclasses aligns naturally with the Template Method and Strategy patterns used in our architecture~\cite{python}.

\subsection{FastAPI}

FastAPI is a modern Python web framework built on Starlette and Pydantic. It was chosen over alternatives such as Flask and Django REST Framework for its combination of high performance (comparable to Node.js), automatic API documentation via Swagger UI, and built-in request validation through Pydantic models~\cite{fastapi}.

\begin{table}[H]
\centering
\caption{Comparison of Python web frameworks}
\label{tab:framework-comparison}
\begin{tabular}{|p{2.5cm}|p{3cm}|p{3cm}|p{3cm}|}
\hline
\textbf{Criteria} & \textbf{FastAPI} & \textbf{Flask} & \textbf{Django REST} \\
\hline
Performance & \cellcolor{green!15}\textbf{High (async)} & Moderate & Moderate \\
\hline
Auto-validation & \cellcolor{green!15}\textbf{Pydantic} & Manual & Serializers \\
\hline
API docs & \cellcolor{green!15}\textbf{Auto (Swagger)} & Manual & Manual \\
\hline
Learning curve & Low & Low & High \\
\hline
Async support & \cellcolor{green!15}\textbf{Native} & Limited & Limited \\
\hline
\end{tabular}
\end{table}

\subsection{Docker and Docker Compose}

Docker~\cite{docker} provides OS-level virtualization through containers, enabling our service to run identically on any machine. Docker Compose~\cite{dockercompose} orchestrates multi-container setups, allowing the entire demonstration environment (policy service + Gitea + Jenkins) to be launched with a single \texttt{docker-compose up} command.

\subsection{Jenkins}

Jenkins~\cite{jenkins} is the most widely adopted open-source automation server. In our project, Jenkins serves as the CI/CD orchestrator that triggers security scans on every code push. Its Jenkinsfile (Pipeline-as-Code) defines the stages: checkout, security scan, build, and deploy. The security scan stage calls our policy service's \texttt{POST /analyze} endpoint and parses the ALLOW/BLOCK verdict.

\subsection{Gitea}

Gitea is a lightweight, self-hosted Git service written in Go. We selected Gitea over GitHub or GitLab for our demonstration because it runs locally without internet access, enabling a fully offline DevSecOps workflow. Gitea provides repository management, webhook triggers, and a web interface at \texttt{localhost:3000}.

\subsection{SQLite}

SQLite is a serverless, zero-configuration relational database engine. In our service, SQLite stores the scan history (timestamps, filenames, violations, decisions) without requiring a separate database server. This choice aligns with our design principle of zero external dependencies.

% ============================================================
\section{Conclusion}
% ============================================================

This chapter established the theoretical and technological foundations upon which our DevSecOps Policy-as-Code microservice is built. We explored the evolution from DevOps to DevSecOps, understanding how the shift-left security paradigm addresses the critical gap between development speed and security rigor. We examined the key concepts of CI/CD, containerization, Infrastructure as Code, and Policy-as-Code that form the conceptual pillars of our project.

The design patterns presented—Layered Architecture, Strategy, Template Method, and Factory—provide the structural framework that enables our service to be extensible, maintainable, and scalable. The technologies selected—Python, FastAPI, Docker, Jenkins, Gitea, and SQLite—were each chosen to fulfill specific requirements while maintaining our commitment to offline capability and simplicity.

With this foundation in place, the following chapter will detail the project planning phase: requirements capture, UML modeling, architecture design, and sprint organization.
